\bookchapter{第十八章：數據備份與遷移}

\chapterintro{吾一鄰人，手機掉入水池中撈出已開不了機矣。其最心疼者非手機本身，乃其中三年未備份過之照片——孩子自出生至會走路之全部記錄。聞其言畢吾心甚感怏怏，因此事完全可避免。此章即認認真真爲數據上一道保險。}

\booksection{一、先搞清：君手機內何物失即無矣}

不同數據之“生存狀態”不同，君須分清：

\begin{itemize}
  \item \textbf{存手機本地、失即真無者}：照片視頻（未開雲同步之話）、備忘錄、短信、通話記錄、下載之文件。
  \item \textbf{存App服務器上、換手機重新登即有者}：微信好友列表、淘寶訂單、銀行App內之交易記錄。此等物存於對應公司之服務器上，非存於君手機內。
  \item \textbf{處於灰色地帶、需尤操作者}：微信聊天記錄（雖微信服務器存了，然君己無法自服務器下載）、手機銀行綁定之設備憑證、兩步驗證App內之密鑰（Google Authenticator那種）。
\end{itemize}

故君當重點護者，乃第一類與第三類。

\booksection{二、整機備份：一次搞定所有物}

整機備份如爲手機拍一完整之“快照”——當前系統設置、桌面佈局、App數據、照片、通訊錄，全部打包。

\booksubsection{方案一：電腦本地備份（最可靠，薦有電腦之人）}

\begin{itemize}
  \item \textbf{iPhone}：將手機以數據線連至電腦上，開iTunes（Windows舊版）或Finder（Mac），於設備頁中選"立即備份"。注意須勾"給備份加密"，否則密碼與健康數據不會被備份入。
  \item \textbf{安卓}：各品牌有己之電腦端軟件，比如華爲手機助手、小米助手、三星Smart Switch。連上電腦按提示操作即行。
  \item 電腦備份之利乃不用花錢買雲空間，且恢復之速比雲端恢復快甚多。弊乃需君記定期去插線備份——大部分人做不到此。
\end{itemize}

\booksubsection{方案二：雲備份（最省心，薦所有人）}

\begin{itemize}
  \item \textbf{iPhone}：設置 君之名 iCloud iCloud雲備份 開之。手機插着電源、連Wi-Fi、鎖屏之時將自動備份。免費之5GB大概率不夠，議花6塊錢升級至50GB。
  \item \textbf{安卓}：各家品牌雲皆有備份之能。華爲手機往設置搜"雲備份"，小米往"小米雲服務"，OPPO/vivo亦類此。又尚有一谷歌自帶之備份（設置 系統 備份），此亦開之。
  \item 雲備份最大之利乃"不用君操心"，自動搞定。最大之弊乃恢復之時須下載甚久，且免費空間一般不夠用。
\end{itemize}

\booksection{三、分類備份之策：要緊者單獨另存一份}

整機備份乃兜底者，然有些物太要緊，值得君單獨另存一份：

\booksubsection{照片與視頻（最珍貴之數字資產）}

\begin{itemize}
  \item \textbf{雙保險方案}：開iCloud照片/Google相冊自動同步，同時每隔數月手將照片複製一份至電腦或移動硬盤上。如此雲端一份、本地一份，互爲備份。
  \item \textbf{不欲花錢之方案}：Google相冊免費予15GB，可設爲"壓縮畫質"以省空間。或每隔一月以數據線將照片全部拷至電腦上。
\end{itemize}

\booksubsection{通訊錄（易忽然極要緊）}

\begin{itemize}
  \item iPhone通訊錄默認與iCloud同步，只要開同步即無需額外備份。然若君欲更安全一點，可往iCloud.com導出vCard文件存電腦上。
  \item 安卓通訊錄通常與谷歌或品牌云同步。往通訊錄App之設置中導出.vcf文件，即一份完整之備份。
\end{itemize}

\booksubsection{微信聊天記錄（單獨處理，重中之重）}

此必單獨拿出言，因整機備份內微信聊天記錄甚多時候不完整，且換機工具根本搬不動其。

\begin{itemize}
  \item \textbf{法一：遷移至新手機}。微信 我 設置 通用 聊天記錄遷移與備份 遷移 遷移至手機/平板。新舊手機連同一Wi-Fi，掃碼即能傳。注意：遷移完成之前，舊手機勿退微信。
  \item \textbf{法二：備份至電腦}。電腦上裝微信PC版，左下角三條槓 遷移與備份 備份與恢復 備份聊天記錄至電腦。後換手機了可再恢復至新手機上。
  \item \textbf{法三：備份至云端}。微信亦支持備份至iCloud或手機雲空間，然注意：此備份只能於君當前手機上恢復，不能跨設備還原。換手機之話，尚須用前二法。
\end{itemize}

\booksection{四、換機指南：自舊手機至新手機}

換新手機乃高興事，然若數據未搬過去即變頭疼事矣。三種情況分別對待：

\booksubsection{同品牌升級（最簡）}

\begin{itemize}
  \item \textbf{iPhone 新iPhone}：開機激活時將舊手機置旁，將自動彈"快速開始"，以舊手機掃新手機上之圖案即行。所有數據、設置、App佈局將原封不動搬過去，包括Wi-Fi密碼皆將自動帶去。大概需二三十分鐘。
  \item \textbf{安卓同品牌}：用品牌自帶之換機工具。華爲用"手機克隆"，小米用"小米換機"，OPPO用"手機搬家"，vivo用"互傳"。步驟皆相仿：新舊手機裝同一App、掃碼連接、勾要搬之物、開傳。
\end{itemize}

\booksubsection{跨品牌換機（安卓換安卓，不同品牌）}

用第三方換機工具或用舊品牌自帶之換機App（大部分品牌之換機App亦支持換至其他安卓手機）。能搬走者有：通訊錄、照片、短信、通話記錄、部分App。搬不走者：App內之登錄狀態（君須重新登）、微信聊天記錄（用微信己之遷移之能）。

\booksubsection{跨平臺換機（安卓換iPhone，或反過來）}

\begin{itemize}
  \item \textbf{安卓 iPhone}：蘋果有一"轉移到iOS"App，安卓手機上下載後，於新iPhone激活時擇"自安卓設備轉移數據"。能搬走通訊錄、照片、短信等基礎數據，然App數據搬不了。微信聊天記錄尚須用微信己之遷移之能。
  \item \textbf{iPhone 安卓}：各安卓品牌有對應之遷移工具，然實言之，跨平臺遷移體驗不如同平臺佳。最要緊之事先單獨做之：通訊錄導出、照片備份、微信聊天記錄遷移。餘者慢慢來。
\end{itemize}

\booksection{五、備份之黃金之則：三二一原則}

此乃數據安全領域之一經典原則，簡而易記：

\begin{itemize}
  \item \textbf{3}：要緊數據至少留三份（一份原件，兩份備份）。
  \item \textbf{2}：備份用二種不同之存儲介質（比如雲端一份，移動硬盤一份）。
  \item \textbf{1}：至少有一份備份置於不同之物理位置（防火災、盜竊等同時毀所有數據）。
\end{itemize}

於常人而言，做到“雲端一份、電腦一份”即已算甚佳矣。無需搞太繁，關鍵乃君真去做了。

\begin{tcbwarn}{微信聊天記錄尤注意}
微信聊天記錄尤注意

\end{tcbwarn}

\begin{tcbpractice}{練之：作一次完整之備份演練}
\begin{enumerate}
  \item 查iCloud雲備份或品牌雲備份有否開啓，確認上一次備份之時。
  \item 查iCloud雲備份或品牌雲備份有否開啓，確認上一次備份之時。
  \item 手備份照片：以數據線連電腦，將DCIM與Pictures文件夾複製至電腦。
  \item 備份微信聊天記錄：開微信 我 設置 通用 聊天記錄遷移與備份 備份至電腦。
  \item 將導出之通訊錄vcf與照片文件夾複製一份至移動硬盤或U盤，成"三二一"內之"一份離線於別處"。
\end{enumerate}

\end{tcbpractice}
